\exercisesection{2}

\begin{enumerate}
\def\labelenumi{\arabic{enumi}.}
\tightlist
\item
  \begin{enumerate}
  \def\labelenumii{(\alph{enumii})}
  \tightlist
  \item
    证明在 Noether 多项式环 \(\mathbf{A} = \mathbf{Z}[\mathbf{X}]\) （\(\mathbf{Z}\) 上一个未定元）中，理想 \(m = 2A + AX\) 是极大的，而理想 \(q = 4A + AX\) 是 m-准素的，但不等于 \(m\) 的任何幂。
  \end{enumerate}
\end{enumerate}

\begin{enumerate}
\def\labelenumi{(\alph{enumi})}
\setcounter{enumi}{1}
\item
  在 Noether 环 \(\mathbf{B} = \mathbf{Z}[2\mathbf{X},\mathbf{X}^2,\mathbf{X}^3] \subset \mathbf{A}\) 中，证明理想 \(\mathfrak{p} = 2\mathrm{BX} + \mathrm{BX}^2\) 是素的，但 \(\mathfrak{p}^2\) 不是 \(\mathfrak{p}\) -准素的，尽管它的根等于 \(\mathfrak{p}\) 。
\item
  设 \(\mathbf{K}\) 是交换域，\(\mathbf{A}\) 是 \(\mathbf{K}[\mathbf{X},\mathbf{Y},\mathbf{Z}]\) 模由 \(Z^2 -\mathbf{XY}\) 生成的理想所得的商环；设 \(x,y,z\) 是 \(\mathbf{X},\mathbf{Y},\mathbf{Z}\) 在 \(\mathbf{A}\) 中的典范像。证明 \(\mathfrak{p} = \mathbf{A}x + \mathbf{A}z\) 是素的，\(\mathfrak{p}^2\) 不是准素的，且 \(\mathfrak{p}^2 = \mathfrak{a}\cap \mathfrak{b}^2\) 是 \(\mathfrak{p}^2\) 的一个准素分解，其中 \(\mathfrak{a} = \mathbf{A}x\) ，\(\mathfrak{b} = \mathbf{A}x + \mathbf{A}y + \mathbf{A}z\) 。
\end{enumerate}

\begin{enumerate}
\def\labelenumi{\arabic{enumi}.}
\setcounter{enumi}{1}
\item
  在 §1 习题 11 (b) 的例子中，证明 \(\mathfrak{a} = \mathfrak{m}^2 \cap \mathfrak{p}_1 \cap \mathfrak{p}_2\) 是 \(\mathfrak{a}\) 在 \(\mathbf{A}\) 中的一个既约准素分解，且 \(\mathfrak{a}\) 关于 \(\mathfrak{m}\) 的饱和等于 \(\mathfrak{a}\) （因而不是准素的）。
\item
  设 A 是 Noether 环，M 是有限生成 A-模。设 Q 是 M 中的 p-准素子模。诸整数 \(n \geqslant 1\) 中使 \(p^{n}M \subset Q\) 者的下确界称为 Q 在 M 中的指数，记作 \(e(M/Q)\) 。
\end{enumerate}

设 \((Q_{\lambda})_{\lambda \in L}\) 是 M 中的一族 \(p\) -准素子模。要使 \(Q = \bigcap_{\lambda \in L} Q_{\lambda}\) 在 M 中是 p-准素的，必须且只需族 \((e(\mathbf{M}/\mathbf{Q}_{\lambda}))_{\lambda\in\mathbf{L}}\) 有上界；此时 \(e(\mathbf{M}/\mathbf{Q})\geqslant e(\mathbf{M}/\mathbf{Q}_{\lambda})\) 对所有 \(\lambda\in L\) 成立。

¶ 4. 设 A 是 Noether 环，M 是有限生成 A-模。设 p 是 Ass(M) 的元素，\(\mathfrak{G}_{\mathfrak{p}}\) 是 M 中满足 \(\operatorname{Ass}(M/Q) = \{\mathfrak{p}\}\) 且 \(\operatorname{Ass}(Q) = \operatorname{Ass}(M) - \{\mathfrak{p}\}\) 的子模 Q 组成的集（由 § 1 第 1 节命题 4，这个集非空）。我们写 \(e_{\mathfrak{p}}(M) = \inf_{\mathbf{Q} \in \mathfrak{G}_{\mathfrak{p}}} e(M/Q)\) （习题 3）。

\begin{enumerate}
\def\labelenumi{(\alph{enumi})}
\item
  设 \(n_{\mathfrak{p}}\) 是整数 \(\geqslant e_{\mathfrak{p}}(\mathbf{M})\) ，\(\mathfrak{F}(n_{\mathfrak{p}})\) 是 \(\mathfrak{G}_{\mathfrak{p}}\) 中由子模 \(\mathbf{Q}\) 之使 \(e(\mathbf{M} / \mathbf{Q}) \leqslant n_{\mathfrak{p}}\) 者组成的子集。证明 \(\mathfrak{F}(n_{\mathfrak{p}})\) 有最小元 \(\mathbf{Q}(\mathfrak{p}, n_{\mathfrak{p}})\) 。
\item
  设 \((n_{\mathfrak{p}})_{\mathfrak{p} \in \operatorname{Ass}(\mathbf{M})}\) 是一族整数，使 \(n_{\mathfrak{p}} \geqslant e_{\mathfrak{p}}(\mathbf{M})\) 对所有 \(\mathfrak{p} \in \operatorname{Ass}(\mathbf{M})\) 成立。证明与这个族相应的诸子模 \(\mathbf{Q}(\mathfrak{p}, n_{\mathfrak{p}})\) 构成 \(\{0\}\) 在 \(\mathbf{M}\) 中的一个由族 \((n_{\mathfrak{p}})\) 典范确定的既约准素分解。若取 \(n_{\mathfrak{p}} = e_{\mathfrak{p}}(\mathbf{M})\) （对所有 \(\mathfrak{p} \in \operatorname{Ass}(\mathbf{M})\) ），则由诸 \(\mathbf{Q}(\mathfrak{p}) = \mathbf{Q}(\mathfrak{p}, e_{\mathfrak{p}}(\mathbf{M}))\) 组成的准素分解称为 \(\{0\}\) 在 \(\mathbf{M}\) 中的典范准素分解。设 \(\{0\} = \bigcap_{\mathfrak{p} \in \operatorname{Ass}(\mathbf{M})} \mathbf{Q}'(\mathfrak{p})\) 是 \(\{0\}\) 在 \(\mathbf{M}\) 中的任意一个既约准素分解。证明 \(e(\mathbf{M}/\mathbf{Q}(\mathfrak{p})) \leqslant e(\mathbf{M}/\mathbf{Q}'(\mathfrak{p}))\) 对所有 \(\mathfrak{p} \in \operatorname{Ass}(\mathbf{M})\) 成立，且若 \(e(\mathbf{M}/\mathbf{Q}(\mathfrak{p})) = e(\mathbf{M}/\mathbf{Q}'(\mathfrak{p}))\) 对某个 \(\mathfrak{p}\) 成立，则 \(\mathbf{Q}(\mathfrak{p}) \subset \mathbf{Q}'(\mathfrak{p})\) （``Ortiz 定理''）。
\item
  证明 \(\mathbf{Q}(\mathfrak{p}, n_{\mathfrak{p}})\) 是 \(\mathfrak{p}^{n_{\mathfrak{p}}} \mathbf{M}\) 关于 \(\mathfrak{p}\) 的饱和（用 § 1 第 2 节命题 6）；\(\mathfrak{p}\) 是 \(\operatorname{Supp}(\mathbf{M} / \mathfrak{p}^n\mathbf{M})\) 的最小元。
\item
  设 S 是 A 的乘性子集，不与素理想 \(\mathfrak{p} \in \operatorname{Ass}(M)\) 相交，Q 是 M 的 p-准素子模。证明 \(e(\mathrm{M}/\mathrm{Q}) = e(\mathrm{S}^{-1}\mathrm{M}/\mathrm{S}^{-1}\mathrm{Q})\) 。证明：若 S 是 A 的乘性子集，\(\Phi\) 是由 \(\mathfrak{p} \in \operatorname{Ass}(M)\) 中使 \(S \cap \mathfrak{p} = \varnothing\) 者组成的集，N 是 \(\{0\}\) 在 M 中关于 S 的饱和，则诸 \(\mathrm{Q}(\mathfrak{p}, n_{\mathfrak{p}})\) （其中 \(\mathfrak{p} \in \Phi\) ）构成 N 在 M 中由族 \((n_{\mathfrak{p}})_{\mathfrak{p} \in \Phi}\) 典范确定的既约准素分解，而诸 \(S^{-1}\mathrm{Q}(\mathfrak{p}, n_{\mathfrak{p}})\) （其中 \(\mathfrak{p} \in \Phi\) ）构成 \(\{0\}\) 在 \(S^{-1}M\) 中由族 \((n_{\mathfrak{p}})_{\mathfrak{p} \in \Phi}\) 典范确定的既约准素分解。考察典范准素分解这一特殊情形。
\end{enumerate}

\begin{enumerate}
\def\labelenumi{\arabic{enumi}.}
\setcounter{enumi}{4}
\item
  设 L 是有限生成的自由 Z-模，T 是阶为素数 p 的幂的有限交换群。证明 \(\{0\}\) 在 \(M = L \oplus T\) 中的典范准素分解（习题 4）是 \(\{0\} = p^{n}L \cap T\) ，其中 n 是最小整数 \(\geqslant 0\) ，使得 \(p^{n}T = 0\) （注意 \(p^{n}M = p^{n}L\) ）。
\item
  在 § 1 习题 5 中，\{0\} 在 A-模 a 中关于 A 的素理想 \{0\} 是准素的。证明 a 的子模 AX 在 a 中关于某个素理想 \(\neq\{0\}\) 是准素的，而 a 的子模 AXY 在 a 中不是准素的。
\item
  设 \(\mathbf{A}\) 是 Noether 环，\(\mathbf{M}\) 是有限生成 A-模，\((\mathfrak{p}_i)_{1\leqslant i\leqslant n}\) 是把 \(\operatorname {Ass}(\mathbf{M})\) 的元素按任意次序排列所得的序列。证明存在一条合成列 \((\mathbf{M}_i)_{0\leqslant i\leqslant n}\) （\(\mathbf{M}\) 的），使得对 \(0\leqslant i\leqslant n - 1\) ，\(\mathbf{M}_{i + 1}\) 是一个 \(p_i\) -准素子模，位于 \(\mathbf{M}_i\) 中（用第 3 节命题 4）。
\end{enumerate}

¶ 8. 设 A 是 Noether 环，E 是有限生成 A-模，F 是 E 的子模，m 是 A 的理想。设 \(\mathbf{F} = \bigcap_{\mathfrak{p} \in \operatorname{Ass}(\mathbf{E}/\mathbf{F})} \mathbf{Q}(\mathfrak{p})\) 是 F 在 E 中的一个既约准素分解。证明 F 在 E 中对于 E 上的 m-进拓扑的闭包等于 \(\bigcap_{\mathfrak{p} \in \Phi} \mathbf{Q}(\mathfrak{p})\) ，其中 \(\Phi\) 是由 \(\mathfrak{p} \in \operatorname{Ass}(\mathbf{E}/\mathbf{F})\) 中使 \(\mathfrak{p} + \mathfrak{m} \neq \mathbf{A}\) 者组成的集。（先考虑 F 在 E 中是 p-准素的情形，并证明若 \(\mathfrak{p} + \mathfrak{m} = \mathbf{A}\) ，则 F 在 E 中稠密；再借助 § 1 习题 15 及第 III 章 § 3 第 4 节定理 3 的推论 1 处理一般情形。）

\begin{enumerate}
\def\labelenumi{\arabic{enumi}.}
\setcounter{enumi}{8}
\item
  设 A 是 Noether 环，m 是 A 的极大理想，M 是使 \(\{0\}\) 在 M 中为 m-准素的 A-模。若 \(\hat{A}\) 是 A 对于 m-进拓扑的 Hausdorff 完备化，证明典范映射 \(M \to M \otimes_{A} \hat{A}\) 是双射（先用第 5 节命题 7 处理 M 有限生成的情形；在一般情形，把 M 看作它的有限生成子模的归纳极限）。
\item
  设 A 是 Noether 环，M 是 A-模。证明下列性质等价：
\end{enumerate}

\((\alpha)\) \(\mathbf{M}\) 的每个有限生成子模都是有限长度的。

(β) M 是有限长度 A-模的归纳极限。

(γ) 每个元素 \(\mathfrak{p} \in \operatorname{Ass}(\mathbf{M})\) 都是 A 的极大理想。

(δ) 每个元素 \(\mathfrak{p} \in \operatorname{Supp}(\mathbf{M})\) 都是 A 的极大理想。

\begin{enumerate}
\def\labelenumi{\arabic{enumi}.}
\setcounter{enumi}{10}
\tightlist
\item
  设 A 是环，M 是 A-模，N 是 M 的子模。使 M/N 上以 a 为比的位似映射为殆幂零的那些 \(a \in A\) 组成的集称为 N 在 M 中的根，记作 \(\mathfrak{r}_{\mathrm{M}}(\mathrm{N})\) 。证明 \(\mathfrak{r}_{\mathrm{M}}(\mathrm{N})\) 是 A 的理想，且对所有 \(p \in \operatorname{Ass}_{f}(\mathrm{M}/\mathrm{N})\) 有 \(\mathfrak{r}_{\mathrm{M}}(\mathrm{N}) \subset p\) 。若 \(N_{1}, N_{2}\) 是 M 的两个子模，则 \(\mathfrak{r}_{\mathrm{M}}(\mathrm{N}_{1} \cap \mathrm{N}_{2}) = \mathfrak{r}_{\mathrm{M}}(\mathrm{N}_{1}) \cap \mathfrak{r}_{\mathrm{M}}(\mathrm{N}_{2})\) ；若 a 是 A 的理想，则 \(\mathfrak{r}_{\mathrm{M}}(\mathfrak{a}\mathrm{N}) \supset \mathfrak{r}(\mathfrak{a}) \cap \mathfrak{r}_{\mathrm{M}}(\mathrm{N})\) ，且 \(\mathfrak{r}_{\mathrm{A}}(\mathfrak{a}) = \mathfrak{r}(\mathfrak{a})\) 。
\end{enumerate}

¶ 12. 设 A 是环，M 是 A-模，N 是 M 的子模。

\begin{enumerate}
\def\labelenumi{(\alph{enumi})}
\item
  要使 \(\operatorname{Ass}_f(\mathbf{M} / \mathbf{N})\) （§ 1，习题 17）只含单个元素 \(\mathfrak{p}\) ，必须且只需 \(\mathbf{N} \neq \mathbf{M}\) ，且对所有 \(a \in \mathbf{A}\) ，以 \(a\) 为比的位似映射在 \(\mathbf{M} / \mathbf{N}\) 上是单射或殆幂零的；此时 \(\mathfrak{p} = \mathfrak{r}_{\mathbf{M}}(\mathbf{N})\) （习题 11），我们也说 \(\mathbf{N}\) 是准素的（或 \(\mathfrak{p}\) -准素的）于 \(\mathbf{M}\) 之中。于是对每个子模 \(\mathbf{M}'\) （\(\mathbf{M}\) 的），只要 \(\mathbf{N} \subset \mathbf{M}'\) 且 \(\mathbf{N} \neq \mathbf{M}'\) ，则 \(\mathbf{N}\) 是 \(\mathfrak{p}\) -准素的于 \(\mathbf{M}'\) 之中。
\item
  若 \(\mathfrak{r}_{\mathbb{M}}(\mathbf{N})\) 是 \(\mathbf{A}\) 的极大理想，证明 \(\mathbf{N}\) 在 \(\mathbf{M}\) 中是准素的。特别地，每个理想 \(q\) （\(\mathbf{A}\) 的），若只含于唯一一个素理想 \(m\) （必为极大），便是 \(m\) -准素的。
\item
  设 \(\mathbf{A}\) 是整环，\(x\) 是 \(\mathbf{A}\) 的元素且 \(\mathfrak{p} = \mathbf{A}x\) 是素的；那么对每个整数 \(n \geqslant 1\) ，\(\mathrm{Ax}^n\) 是 \(\mathfrak{p}\) -准素的。证明这个结果不能推广到 A 不是整环的情形（取 A = B/b，其中 B = K{[}X, Y{]} （K 是域），b = BX² + BXY）。
\item
  若 A 是绝对平坦环（第 I 章，§ 2，习题 17），则 A 的准素理想与 A 的素理想相同。
\item
  给出一个 Z-模 M 的例子，使 \(\{0\}\) 是 p-准素的（p 为素数），但对 Z 中任何 \(a \neq 0\) ，位似映射 \(a_{M}\) 都不是幂零的（参见第 II 章，§2，习题 3）。
\end{enumerate}

\begin{enumerate}
\def\labelenumi{\arabic{enumi}.}
\setcounter{enumi}{12}
\tightlist
\item
  \begin{enumerate}
  \def\labelenumii{(\alph{enumii})}
  \tightlist
  \item
    设 \(u: \mathbf{M} \to \mathbf{M}'\) 是满的 A-模同态。证明：若 \(\mathbf{N}'\) 是 \(\mathbf{M}'\) 中的准素子模，则 \(\mathbf{N} = \bar{u}^{-1}(\mathbf{N}')\) 在 \(\mathbf{M}\) 中是准素的，且 \(\mathfrak{r}_{\mathbf{M}}(\mathbf{N}) = \mathfrak{r}_{\mathbf{M}'}(\mathbf{N}')\) 。
  \end{enumerate}
\end{enumerate}

\begin{enumerate}
\def\labelenumi{(\alph{enumi})}
\setcounter{enumi}{1}
\item
  A-模 M 中 p-准素子模组成的非空有限族的交仍是 M 中的 p-准素子模。该命题能否推广到任意交？
\item
  要使 A-模 M 满足 \(\{0\}\) 在 M 中是 p-准素的，必须且只需对 M 的每个子模 \(N \neq 0\) 有 \(\operatorname{Supp}(N) = V(p)\) 。（注意对所有 \(x \in M\) ，M 含有一个同构于 \(A/\operatorname{Ann}(x)\) 的子模。）
\end{enumerate}

\begin{enumerate}
\def\labelenumi{\arabic{enumi}.}
\setcounter{enumi}{13}
\tightlist
\item
  \begin{enumerate}
  \def\labelenumii{(\alph{enumii})}
  \tightlist
  \item
    把第 1 节的命题 3 推广到任意环。
  \end{enumerate}
\end{enumerate}

\begin{enumerate}
\def\labelenumi{(\alph{enumi})}
\setcounter{enumi}{1}
\tightlist
\item
  设 M 是 A-模，使得对 A 的每个乘性子集 S，典范映射 \(M \rightarrow S^{-1}M\) 的核为 \(\{0\}\) 或 M。证明 \(\{0\}\) 在 M 中是准素的（对 A 中所有 \(a \neq 0\) ，考虑由 \(a^{n}\) （ \(n \geqslant 0\) ）组成的乘性子集）。
\end{enumerate}

¶ 15. 设 q 是环 A 中的准素理想，p 是它的根。证明在多项式环 B = A{[}X{]} 中，理想 Bq 是准素的，且其根为理想 Bp。（设 \(f(X) = \sum_{j} a_j X^j\) 与 \(g(X) = \sum_{j} b_j X^j\) 是两个非常数多项式，满足 \(fg \in Bq\) 且 \(f \notin Bp\) 。设 \(a_m\) 是不属于 p 的最小下标系数；设 a 是由 \(a_0, a_1, \ldots, a_{m-1}\) 生成的 A 的理想；则存在整数 k 使 \(a^k \subset q\) 。设 \(q_i\) 是 \(a^{k-i}\) 在 q 中的传递子（\(i \leqslant k\)）。对 i 作归纳证明 \(g \in Bq_i\) ，用反证法论证。

\begin{enumerate}
\def\labelenumi{\arabic{enumi}.}
\setcounter{enumi}{15}
\item
  在环 A 中，设 p 是非极大的素理想，q 是异于 p 的 p-准素理想。设 x 是 p - q 的元素，y 是 A - p 的元素且 \(p + Ay \neq A\) 。证明理想 \(a = q + Axy\) （它满足 \(p \supset a \supset q\) ）不是准素的（注意 \(x \in a\) 是不可能的）。
\item
  设 M 是 A-模，N 是 M 中的 p-准素子模。
\end{enumerate}

\begin{enumerate}
\def\labelenumi{(\alph{enumi})}
\item
  设 F 是 M 的不含于 N 的子模。证明传递子 N:F 是 A 中含于 p 的理想。
\item
  设 \(\mathfrak{a}\) 是 \(\mathbf{A}\) 的理想。证明：若 \(\mathfrak{a} \subset \mathbb{N}: \mathbb{M}\) ，则 \(\mathbb{N}: a = \mathbb{M}\) ；若 \(\mathfrak{a} \notin \mathbb{N}: \mathbb{M}\) ，则 \(\mathbb{N}: \mathfrak{a}\) 是一个 \(\mathfrak{p}\) -准素子模，位于 \(\mathbb{M}\) 中；若 \(\mathfrak{a} \notin \mathfrak{p}\) ，则 \(\mathbb{N}: \mathfrak{a} = \mathbb{N}\) 。
\end{enumerate}

\begin{enumerate}
\def\labelenumi{\arabic{enumi}.}
\setcounter{enumi}{17}
\tightlist
\item
  \begin{enumerate}
  \def\labelenumii{(\alph{enumii})}
  \tightlist
  \item
    设 \(\mathbf{A}\) 是环，\(\mathbf{M}\) 是 \(\mathbf{A}\) -模。证明：若 \(\mathfrak{p}\) 是 \(\operatorname{Ass}_f(\mathbf{M})\) 的极小元，则 \(\{0\}\) 关于 \(\mathfrak{p}\) 在 \(\mathbf{M}\) 中的饱和是 \(\mathfrak{p}\) -准素的于 \(\mathbf{M}\) 之中（参见 § 1，习题 17 (b)）。
  \end{enumerate}
\end{enumerate}

\begin{enumerate}
\def\labelenumi{(\alph{enumi})}
\setcounter{enumi}{1}
\item
  设 p 是 A 的素理想。证明对每个整数 n \textgreater{} 0 ，\(p^{n}\) 在 A 中关于 p 的饱和在 A 中是 p-准素的（利用习题 14 (a)）；该饱和记作 \(p^{(n)}\) ，称为 p 的 n 次符号幂。
\item
  证明：若 A 的每个素理想都是极大的，则对每个 A-模 M 及 M 的每个子模 N，N 是 M 中一个准素子模族（有限或无限）的交（利用 (a) 及 § 1 的习题 17 (i)）。
\end{enumerate}

* 19. 确定高为 2 的赋值的环中的准素理想（参见第 VI 章）；由此给出一个环的例子，其中存在不能写成准素理想族（有限或无限）的交的理想，尽管该环中只有两个不同的素理想。 *

¶ 20. 准素分解与既约准素分解的概念可以对任意环 A 像第 2 节与第 3 节中那样定义（利用习题 12 中定义的准素子模概念）。

\begin{enumerate}
\def\labelenumi{(\alph{enumi})}
\item
  把命题 4 与命题 6 推广：处处把 Ass 换成 \(\mathbf{Ass}_f\) 。
\item
  设 E 是有限生成 A-模，F 是 E 的子模。对 A 的每个乘性子集 S，用 \(\operatorname{sat}_{\mathbf{S}}(\mathbf{F})\) 表示 F 在 E 中关于 S 的饱和。证明：若 F 在 E 中有一个准素分解，则 F 在 E 中的每个饱和都形如 F: (a) ，其中某个 \(a \in A\) 。（设 \(p_{i}\) （ \(1 \leqslant i \leqslant r\) ）是 \(\operatorname{Ass}_{f}(\operatorname{E}/\operatorname{F})\) 的元素。先利用 (a) 证明，对 A 的每个乘性子集 S，存在 \(b \in A\) ，使得若 T 是由 \(b^{n}\) （ \(n \geqslant 0\) ）组成的乘性子集，则 \(\operatorname{sat}_{\mathbf{T}}(\mathbf{F}) = \operatorname{sat}_{\mathbf{S}}(\mathbf{F})\) ，其中 b 这样选取：\(b \in p_{i}\) 当 \(p_{i} \cap S \neq 0\) 时，而 \(b \notin p_{i}\) 否则。然后证明当 m 充分大时可取 \(a = b^{m}\) 。）
\item
  给出一个 \(\mathbf{Z}\) -模 E 的例子，使得 \(\{0\}\) 关于 E 是准素的，但存在不形如 0: (a) 的饱和 \(\mathrm{sat}_{\mathbb{S}}(0)\) （参见《代数》，第 VII 章，§ 2，习题 3）。
\end{enumerate}

\begin{enumerate}
\def\labelenumi{\arabic{enumi}.}
\setcounter{enumi}{20}
\tightlist
\item
  设 A 是环，F 是 A{[}X{]} 的首一多项式，B 是环 A{[}X{]}/F.A{[}X{]} ，m 是 A 的极大理想，k = A/m 是剩余域，\(\bar{F}\) 是 F 在 k{[}X{]} 中的典范像。
\end{enumerate}

\begin{enumerate}
\def\labelenumi{(\alph{enumi})}
\item
  设 \(\overline{\mathbf{F}} = \prod_{i=1}^{n} f_i^{e_i}\) ，其中 \(f_i\) 是 \(k[\mathbf{X}]\) 中互不相同的不可约多项式。设 \(\mathbf{F}_i \in \mathbf{A}[\mathbf{X}]\) 是首一多项式，使得它在 \(k[\mathbf{X}]\) 中的典范像是 \(f_i\) 。设 \(\mathfrak{M}_i\) 表示理想 \(\mathsf{mB} + \mathbf{F}_i\) . B（B 的）。证明诸 \(\mathfrak{M}_i\) 互不相同，且是 B 中含 \(\mathsf{mB}\) 的仅有的极大理想（注意 \(\mathbf{B}/\mathsf{mB}\) 是 \(\mathbf{A}/\mathsf{m} = k\) 上由 \(\overline{\mathbf{F}}\) 的根生成的有限秩代数）。
\item
  对所有 i，记 \(Q_{i} = mB + F_{i}^{e_{i}}\) . B。证明诸 \(Q_{i}\) 在 B 中是 \(M_{i}\) -准素的，且 \(mB = Q_{1} \cap Q_{2} \cap \cdots \cap Q_{n} = Q_{1}Q_{2} \cdots Q_{n}\) 。
\item
  要使 B 是局部环，必须且只需 A 是局部环，且 \(\overline{F}\) 是 k{[}X{]} 中某个不可约多项式的幂。
\end{enumerate}

¶ 22. 设 E 是 A-模，F 是 E 的子模。

\begin{enumerate}
\def\labelenumi{(\alph{enumi})}
\item
  设 \(\mathfrak{p}\) 是 \(\mathbf{A}\) 的素理想，\(a\) 是 \(\mathbf{A}\) 的元素，使得 \(a \notin \mathfrak{p}\) ，且 \(\mathbf{Q} = \mathbf{F}\) : (a) 是 \(\mathfrak{p}\) -准素的于 \(\mathbf{E}\) 之中。证明 \(\mathbf{F} = \mathbf{Q} \cap (\mathbf{F} + a\mathbf{E})\) 。
\item
  假设存在 \(b \in \mathbf{A}\) 与 \(x \in \mathbf{E}\) ，使得 \(b^n x \notin \mathbf{F}\) 对所有 \(n > 0\) 成立。证明：若存在整数 \(n > 0\) 使 \(\mathbf{F}: b^{n+1} = \mathbf{F}: b^n\) ，则 \(\mathbf{F} = (\mathbf{F} + \mathbf{Ab}^n x) \cap (\mathbf{F}: (b^n))\) 。
\item
  假设 \(\mathbf{F}\) 关于 \(\mathbf{E}\) 是不可约的（《代数》，第 II 章，§ 2，习题 16），并考察以下三个性质：
\end{enumerate}

(α) F 在 E 中是准素的。

(β) 对 A 的每个素理想 p，F 在 E 中关于 p 的饱和都形如 F: (a) ，其中某个 \(a \notin p\) 。

\((\gamma)\) 对所有 \(b \in \mathbf{A}\) ，序列 \((\mathbf{F}:(b^n))_{n \geqslant 1}\) 是平稳的。

证明条件 \((\beta)\) 、 \((\gamma)\) 中的每一个都蕴涵 \((\alpha)\) （用 (a)、(b) 及习题 18 (a)）。若 E 是有限生成的，则三个条件 \((\alpha)\) 、 \((\beta)\) 与 \((\gamma)\) 等价（参见习题 20 (b) 与 20 (c)）。

\begin{enumerate}
\def\labelenumi{(\alph{enumi})}
\setcounter{enumi}{3}
\item
  设 \(\mathbf{K}\) 是域，\(\mathbf{A}\) 是环 \(\mathbf{K}[\mathbf{X},\mathbf{Y}]\) ，\(\mathfrak{m}\) 是极大理想 \(\mathbf{AX} + \mathbf{AY}\) （属于 \(\mathbf{A}\) ）；证明理想 \(\mathfrak{m}^2\) 是准素的，但不是不可约的。
\item
  设 A 是不必交换的环，E 是左 Noether A-模。证明 E 的每个子模 F 都是 E 中一个有限不可约子模族的交（考虑 E 的一个子模，它在那些不能写成不可约子模的有限交的子模之中是极大的）。
\item
  由 (e) 与 (c) 推出第 2 节定理 1 的一个新证明。
\end{enumerate}

¶ 23. 设 A 是环。若 A-模 E 是有限生成的，且 E 的每个子模在 E 中都有一个准素分解，则称 E 为 Lasker 模。

\begin{enumerate}
\def\labelenumi{(\alph{enumi})}
\tightlist
\item
  要使有限生成 A-模 E 是 Lasker 模，必须且只需它满足下面两条公理：
\end{enumerate}

\((\mathbf{LA}_{\mathbf{I}})\) 对每个子模 \(\mathbf{F}\) （\(\mathbf{E}\) 的）及每个素理想 \(\mathfrak{p}\) （\(\mathbf{A}\) 的），\(\mathbf{F}\) 关于 \(\mathfrak{p}\) 在 \(\mathbf{E}\) 中的饱和都形如 \(\mathbf{F}:(a)\) ，其中某个 \(a\notin\mathfrak{p}\) 。

\((\mathbf{L}\mathbf{A}_{\mathrm{II}})\) 对每个子模 \(\mathbf{F}\) （\(\mathbf{E}\) 的），每个递减序列 \((\mathrm{sat}_{\mathbf{S}_{n}}(\mathbf{F}))\) （其中 \((\mathbf{S}_{n})\) 是 \(\mathbf{A}\) 的乘性子集的任意递减序列）都是平稳的。

（为证条件充分，先用 \((\mathrm{LA}_{\mathrm{I}})\) 及习题 18 (a) 与 22 (a) 证明：对 E 的每个子模 F，存在 E 的子模 Q，它关于某个理想 \(p \supset F: E\) 是准素的，且存在子模 \(G = F + aE\) ，其中 \(a \notin p\) ，使得 \(F = Q \cap G\) 。然后用反证法论证：证明这时将存在无穷序列 \((\mathbf{Q}_{n})\) （由 E 的 \(p_{n}\) -准素子模组成）与 E 的子模的严格递增序列 \((\mathbf{G}_{n})\) ，使得若记 \(H_{n} = Q_{1} \cap Q_{2} \cap \cdots \cap Q_{n}\) ，则：(1) \(F = H_{n} \cap G_{n}\) ；(2) \(G_{n}\) 在含 \(G_{n-1}\) 且使 \(\mathbf{F} = \mathbf{H}_n \cap \mathbf{G}\) 的诸子模 G 之中是极大的；(3) 存在 \(a_n \notin p_n\) 使 \(a_n \mathrm{E} \subset \mathrm{G}_n\) 。再证明，若 \(\mathrm{S}_n = \bigcap_{1 \leqslant i \leqslant n} (\mathrm{A} - p_i)\) ，则 \(\mathrm{S}_n\) 与 \(\mathrm{G}_n: \mathrm{E}\) 相交，并由此推出 \(\operatorname{sat}_{\mathrm{S}_n}(\mathrm{F}) = \mathrm{H}_n\) 。最后借助 \((\mathrm{LA}_{\mathrm{II}})\) 完成证明。）

* (b) 给出一个环 A 的例子，使得 A-模 A 不满足 \((\mathrm{LA}_{\mathrm{I}})\) ，但其中每个理想 \(a \subset A\) 都不可约，且 A 的所有乘性子集 S 的饱和 \(\operatorname{sat}_{\mathrm{S}}(a)\) 组成的集是有限的（参见习题 19）。 *

\begin{enumerate}
\def\labelenumi{(\alph{enumi})}
\setcounter{enumi}{2}
\item
  设 K 是域，\(B = K[[X]]\) 是 K 上的形式幂级数代数，m 是 B 中由无常数项的形式幂级数组成的极大理想。在乘积代数 \(B^{N}\) 中，考虑由 1 与理想 \(n = m^{(N)}\) 生成的子代数 A。证明在 A 中，异于 n 的素理想只有那样的理想，它们是 n 中除一个分量理想外所有分量理想的直和；于是 A 的素理想的严格递增序列至多含两个元素。设 \(a \neq A\) 是 A 的理想，S 是 A 的不与 a 相交的乘性子集；若 \(a_{n}\) （相应地 \(S_{n}\) ）是 a（相应地 S）在 \(B_{n}\) （\(B^{N}\) 的第 n 个因子）上的投影，证明 \(\text{sat}_{S}(a)\) 是诸理想 \(\text{sat}_{S_{n}}(a_{n})\) 的直和（注意，若 \(s \in S\) ，\(s_{n} \in S_{n}\) 是 s 的第 n 个投影，且 \(x_{n} \in \text{sat}_{S_{n}}(a_{n})\) 使得 \(s_{n}x_{n} \in a_{n}\) ，则 \(s^{2}x_{n} \in a\) ）。由此推出 A 满足公理 \((\mathbf{LA}_{\mathbf{I}})\) 但不满足 \((\mathbf{LA}_{\mathbf{II}})\) （用习题 20 (b)）。
\item
  证明：若 A-模 E 满足公理 \((\mathrm{LA}_{\mathrm{I}})\) ，则 E 的每个子模都是 E 中一个准素子模族（有限或无限）的交（像 (a) 中那样论证，用超限归纳法）。
\item
  设 E 是有限生成 A-模，\(E' \subset E\) 是有限生成子模。要使 E 是 Lasker 模，必须且只需 \(E'\) 与 \(E/E'\) 都是 Lasker 模。特别地，若 A 是 Lasker 环，则每个有限生成 A-模都是 Lasker 模。若 E 是 Lasker A-模，则对 A 的每个乘性子集 S，\(S^{-1}E\) 是 Lasker \(S^{-1}A\) -模。
\end{enumerate}

¶ 24. 在 Lasker 环 A（习题 23）中，设 \(p_1, p_2, p_3\) 是三个不同的素理想，满足 \(p_1 \subset p_2 \subset p_3\) ；设 \(x_2 \in p_2 - p_1, x_3 \in p_3 - p_2\) 。如有必要，把 \(p_2\) 与 \(p_3\) 分别换成含于 \(p_2, p_3\) 且分别含 \(x_2\) 与 \(x_3\) 的素理想，我们便可设 \(p_2\) 在 \(\operatorname{Ass}_f(A / (p_1 + Ax_2))\) 中极小，且 \(p_3\) 在 \(\operatorname{Ass}_f(A / p_2 + Ax_3)\) 中极小。考察理想 \(a = p_1 + x_2p_3\) 。

\begin{enumerate}
\def\labelenumi{(\alph{enumi})}
\item
  证明 \(x_{2} \notin \mathfrak{a}\) ，且 \(x_{3}^{k} \notin \mathfrak{a}\) 对每个整数 \(k > 0\) 成立；由此推出 \(\mathfrak{a}\) 在 A 中不是准素的。证明 \(\mathfrak{p}_{2}\) 是 \(\operatorname{Ass}_f(\mathbf{A} / \mathfrak{a})\) 的极小元。
\item
  设 \(\mathfrak{a} = \bigcap_{i=1}^{n} \mathfrak{q}_i\) 是 \(\mathfrak{a}\) 在 \(\mathbf{A}\) 中的一个既约准素分解，\(\mathfrak{p}_i'\) 是 \(\mathfrak{q}_i\) 的根；设 \(\mathfrak{p}_3 \notin \mathfrak{p}_i'\) 对 \(1 \leqslant i \leqslant s\) 成立，\(\mathfrak{p}_3 \subset \mathfrak{p}_i'\) 对 \(s + 1 \leqslant i \leqslant n\) 成立；证明必有 \(s < n\) ，且 \(x_2 \in \mathfrak{q}_i\) 对 \(1 \leqslant i \leqslant s\) 成立。证明存在指标 \(i \geqslant s + 1\) 使 \(\mathfrak{p}_3 = \mathfrak{p}_i'\) 。（用反证法论证：在相反情形，将存在 \(y \notin \mathfrak{p}_3\) 使得 \(y \in \mathfrak{q}_i\) 对所有 \(i \geqslant s + 1\) 成立；于是 \(x_2y \in \mathfrak{a}\) ，证明这个关系与 \(y \notin \mathfrak{p}_3\) 矛盾。）
\end{enumerate}

由此得出 \(\mathfrak{p}_3\) 是 \(\operatorname{Ass}_f(A / a)\) 的一个非极小元；我们用 \(\mathfrak{q}_3'\) 表示 \(a\) 在 \(A\) 中的、在上面的分解中对应于它的准素分量。

\begin{enumerate}
\def\labelenumi{(\alph{enumi})}
\setcounter{enumi}{2}
\tightlist
\item
  设 \(\mathfrak{b} = \mathfrak{p}_2 \cap \mathfrak{q}_3'\) ；另一方面，设 \(m > 0\) 使得 \(x_3^m \in \mathfrak{q}_3'\) ，并设 \(\mathfrak{c} = \mathfrak{b} + Ax_3^{2m}\) 。证明 \(\mathfrak{c}\) 关于 \(\mathfrak{p}_3\) 的饱和是一个 \(\mathfrak{p}_3\) -准素理想 \(\mathfrak{q}_3''\) （证明 \(\mathfrak{p}_3\) 的每个元素都有一个幂在 \(\mathfrak{q}_3''\) 中，方法是利用 \(\mathfrak{p}_2 + Ax_3\) 关于 \(\mathfrak{p}_3\) 的饱和是准素的这一事实）。另一方面证明 \(x_3^m \notin \mathfrak{q}_3''\) ，且 \(\mathfrak{b} = \mathfrak{p}_2 \cap \mathfrak{q}_3''\) ；于是理想 \(\mathfrak{b}\) 有两个不同的既约准素分解。
\end{enumerate}

\begin{enumerate}
\def\labelenumi{\arabic{enumi}.}
\setcounter{enumi}{24}
\tightlist
\item
  设 A 是 Lasker 环。设 \(a\) 是 A 的理想，\(a = \bigcap_{1 \leqslant i \leqslant n} q_i\) 是 \(a\) 在 A 中的一个既约准素分解；用 \(p_i\) 表示 \(q_i\) 的根；设 \(\operatorname{Ass}_f(A/a)\) 的极小元是诸 \(p_i\) （ \(1 \leqslant i \leqslant s\) ），且 \(s < n\) 。记 \(b = \bigcap_{1 \leqslant i \leqslant s} q_i\) ，\(c = \bigcap_{s+1 \leqslant i \leqslant n} q_i\) 。
\end{enumerate}

\begin{enumerate}
\def\labelenumi{(\alph{enumi})}
\item
  证明存在 \(x \in c\) 使得 \(x \notin p_i\) 对 \(1 \leqslant i \leqslant s\) 成立，且 \(a = b \cap (a + Ax)\) 。
\item
  例如设 \(p_{s+1}\) 在 \(\operatorname{Ass}_{f}(\mathbf{A}/c)\) 中极小。证明 \(xp_{s+1} + a \neq Ax + a\) ，且 \(a = b \cap (xp_{s+1} + a)\) ；此外，\(p_{s+1}\) 在 \(\operatorname{Ass}_{f}(\mathbf{A}/b)\) 中极小，其中 \(d = xp_{s+1} + a\) 。证明 d 关于 \(p_{s+1}\) 的饱和异于 \(q_{s+1}\) 。（注意这个饱和不可能含 x。）
\end{enumerate}

¶ 26. 设 A 是一个 Lasker 环，其中每个异于 A 的理想在 A 中都有唯一的既约准素分解。

\begin{enumerate}
\def\labelenumi{(\alph{enumi})}
\item
  证明对每个理想 \(\mathfrak{a} \neq \mathbf{A}\) ，\(\operatorname{Ass}_f(\mathbf{A} / \mathfrak{a})\) 的所有元素在这个集中都是极小的（用习题 25 (b)）。设 \(\{0\} = \bigcap_{1 \leqslant i \leqslant n} q_i\) 是 \(\mathbf{A}\) 中的一个既约准素分解，\(p_i\) 是 \(q_i\) （ \(1 \leqslant i \leqslant n\) ）的根。证明 \(\mathbf{A}\) 的每个异于诸 \(p_i\) 的素理想都是极大的（用习题 24 (c)）。若 \(p\) 是 \(\mathbf{A}\) 的素理想，\(q\) 是含于 \(p\) 且使 \(p / q\) 在 \(\mathbf{A} / q\) 中为幂零理想的理想，证明每个理想 \(\mathfrak{a}\) （满足 \(q \subset \mathfrak{a} \subset p\) ）都是 \(p\) -准素的。由此推出，若某个 \(p_i\) 不是极大的，则 \(p_i\) 是唯一的 \(p_i\) -准素理想，且 \(p_i^2 = p_i\) （用习题 16）。
\item
  证明在 A 中，每个由等于自身根的理想组成的递增序列都是平稳的。由此推出，每个严格递增的理想序列 \((\mathfrak{a}_n)\) ，若每个商 \(\mathfrak{a}_n / \mathfrak{a}_{n-1}\) 在 \(\mathbf{A} / \mathfrak{a}_{n-1}\) 中含非幂零元素，则它是有限的。由此得出，对 A 的每个素理想 \(\mathfrak{p}\) ，存在一个以 \(\mathfrak{p}\) 为根的准素理想，它是有限生成的。特别地，诸 \(\mathfrak{p}_i\) 中每个非极大的都是有限生成的，从而含有一个生成它的幂等元 \(e_i\) （第 II 章，§ 4，习题 15）。
\item
  证明对 \(i \neq j \mathfrak{p}_i + \mathfrak{p}_j = \mathbf{A}\) （考虑 \(e_i + e_j - e_i e_j\) ，若 \(\mathfrak{p}_i\) 与 \(\mathfrak{p}_j\) 不是极大的）。由此得出 \(\mathbf{A}\) 同构于有限多个环 \(\mathbf{A}_i = \mathbf{A} / \mathfrak{q}_i\) 的乘积，这些环满足：或者 \(\mathbf{A}_i\) 是局部环，其极大理想是幂零理想；或者 \(\mathbf{A}\) 是一个整环，其中每个素理想 \(\neq \{0\}\) 都是极大的，且每个理想 \(\neq\{0\}\) 只含于有限多个极大理想 * （参见 § 1，习题 2）。* 证明逆命题。
\end{enumerate}

¶ 27. 设 M 是 A-模，Q 是 M 的 p-准素子模。若 m 是使 \(p^{k}M \subset Q\) 的最小整数 k（参见习题 3），则称 Q 的指数为 m；若不存在具有这一性质的整数 k，则称 Q 的指数为无穷。若 Q 的指数有限，也称 Q 是强准素的。

\begin{enumerate}
\def\labelenumi{(\alph{enumi})}
\tightlist
\item
  证明：若 \(\mathbf{Q} \colon \mathfrak{p} = \mathbf{Q}\) ，则 \(\mathbf{Q}\) 的指数为无穷，且存在 \(\alpha \in \mathfrak{p}\) 使 \(\mathbf{Q} \colon (\alpha)\) 异于 \(\mathbf{Q}\) ，是 \(\mathfrak{p}\) -准素子模于 \(\mathbf{M}\) 之中，且指数为无穷。
\end{enumerate}

* (b) 给出一个环 A 及一个指数无穷的 p-准素理想 q 的例子，使得 q: p = q（参见 § 1，习题 2）。 *

\begin{enumerate}
\def\labelenumi{(\alph{enumi})}
\setcounter{enumi}{2}
\item
  设 K 是域，A 是多项式环 \(K[X_{n}]_{n \in N}\) ，p 是 A 中由诸 \(X_{n}\) 生成的素理想，q 是由乘积 \(X_{i}X_{j} (i \neq j)\) 与元素 \(X_{n}^{n+1}\) 生成的 p-准素理想。证明对所有 \(\alpha \in p - q\) ，\(q: (\alpha)\) 的指数有限，\(q: p \neq q\) ，且 q: p 是指数无穷的 p-准素理想。
\item
  设 \(\mathbf{Q}\) 是 \(p\) -准素子模（在 \(\mathbf{M}\) 中），指数有限，为 \(e\) 。证明诸 \(p^k\mathbf{M}\) （ \(k < e\) ）互不相同，且诸子模 \(\mathbf{Q}: p^k\) （ \(k < e\) ）互不相同。对每个理想 \(a\) （\(A\) 中的），证明或者 \(Q: a = Q\) ，或者 \(Q: a\) 是一个 \(p\) -准素子模，位于 \(M\) 中，异于 \(M\) ，且指数 \(\leqslant e - 1\) （注意若 \(a \subset p\) 则 \(p^{e-1}\mathbf{M} \subset Q: a\) ）。
\end{enumerate}

¶ 28. 若有限生成 A-模 E 的每个子模在 E 中都有一个其元素都是强准素的准素分解（习题 27），则称 E 为强 Lasker 模。

\begin{enumerate}
\def\labelenumi{(\alph{enumi})}
\tightlist
\item
  证明：要使有限生成 A-模 E 是强 Lasker 模，必须且只需它满足习题 23 的公理 (LA \(_{II}\) ) 及下面的公理：
\end{enumerate}

\((\mathbf{L}\mathbf{A}_{\mathrm{III}})\) 对理想的每个序列 \((\mathfrak{b}_k)\) （A 中的）及 E 的每个子模 F，子模的递增序列 \(\mathbf{F}:(\mathfrak{b}_1\mathfrak{b}_2\dots \mathfrak{b}_k)\) 是平稳的。

（为证条件必要，用习题 27 (d)。为证条件充分，先证明 \((\mathrm{LA}_{\mathrm{III}})\) 蕴涵 \((\mathrm{LA}_{\mathrm{I}})\) ，再利用习题 27 (a) 证明它蕴涵每个准素理想的指数都有限。）

\begin{enumerate}
\def\labelenumi{(\alph{enumi})}
\setcounter{enumi}{1}
\item
  证明习题 23 (c) 中定义的环 A 满足 \((\mathbf{LA}_{\mathrm{III}})\) 。
\item
  设 \(\mathbf{K}\) 是域，\(\mathbf{V}\) 是无穷维 \(\mathbf{K}\) -向量空间，\(\mathbf{A} = \mathbf{K} \oplus \mathbf{V}\) 是一个环，其乘法由下式定义
\end{enumerate}

\[
(a, x) (a ^ {\prime}, x ^ {\prime}) = (a a ^ {\prime}, a x ^ {\prime} + a ^ {\prime} x).
\]

证明在 A 中每个理想 \(\neq A\) 都是强准素的，尽管 A 不是 Noether 环。

\begin{enumerate}
\def\labelenumi{(\alph{enumi})}
\setcounter{enumi}{3}
\item
  设 E 是有限生成 A-模，\(E' \subset E\) 是有限生成子模。要使 E 是强 Lasker 模，必须且只需 \(E'\) 与 \(E/E'\) 都是强 Lasker 模。特别地，若 A 是强 Lasker 环，则每个有限生成 A-模都是强 Lasker 模。若 E 是强 Lasker A-模，则对 A 的每个乘性子集 S，\(S^{-1}E\) 是强 Lasker \(S^{-1}A\) -模。
\item
  把习题 3 与 4 推广到强 Lasker 模。
\end{enumerate}

¶ 29. 设 A 是强 Lasker 环（习题 28）。

\begin{enumerate}
\def\labelenumi{(\alph{enumi})}
\item
  设 \(\mathfrak{a}, \mathfrak{b}\) 是 \(\mathbf{A}\) 的两个理想，\(\mathfrak{ab} = \bigcap_{i=1}^{n} q_i\) 是 \(\mathfrak{ab}\) 在 \(\mathbf{A}\) 中的一个既约准素分解，\(\mathfrak{c}\) 是诸 \(q_i\) （其根不含 \(\mathfrak{b}\) 者）的交；证明 \(\mathfrak{a} \subset \mathfrak{c}\) （若 \(q_i\) 的根 \(p_i\) 不含 \(\mathfrak{b}\) ，且 \(x \in \mathfrak{b}\) 而 \(x \notin p_i\) ，考虑乘积 \(ax\) ）。
\item
  由 (a) 推出：存在三个理想 \(\mathfrak{a}', \mathfrak{b}', \mathfrak{c}\) 及一个整数 \(s > 0\) ，使得 \(\mathfrak{a}^s \subset \mathfrak{a}', \mathfrak{b}^s \subset \mathfrak{b}', (\mathfrak{a} + \mathfrak{b})^s \subset \mathfrak{c}\) ，且
\end{enumerate}

\[
a b = a ^ {\prime} \cap b = a \cap b ^ {\prime} = a \cap b \cap c.
\]

\begin{enumerate}
\def\labelenumi{(\alph{enumi})}
\setcounter{enumi}{2}
\item
  证明对 A 的理想的每个有限族 \((\mathfrak{a}_k)\) ，存在整数 \(m > 0\) 使 \(\bigcap_{k} \mathfrak{a}_k^m \subset \prod_{k} \mathfrak{a}_k\) （应用 (b)，对理想 \(\mathfrak{a}_k\) 的个数作归纳论证）。
\item
  对每个理想 \(\mathfrak{a}\) （\(\mathbf{A}\) 的），证明交 \(\mathfrak{b} = \bigcap_{n=1} \mathfrak{a}^n\) 是由那些 \(x \in \mathbf{A}\) （满足 \(x \in x\mathfrak{a}\) ）组成的理想（把 (b) 应用于乘积 \((\mathbf{A}x)\mathfrak{a}\) ）。由此推出 \(\mathfrak{b}\) 是 \(\{0\}\) （对 \(\{0\}\) 的一个既约准素分解而言）的其根与 \(1 + \mathfrak{a}\) 相交的那些准素分量的交。
\item
  由 (d) 推出，对每个素理想 \(\mathfrak{p}\) （\(\mathbf{A}\) 的），符号幂 \(\mathfrak{p}^{(n)}\) （习题 18 (b)）对 \(n\in\mathbf{N}\) 的交是由那些 \(x\in\mathbf{A}\) （满足 \(x\in x\mathfrak{p}\) ）组成的理想（考虑环 \(\mathbf{A}_{\mathfrak{p}}\) ）。
\end{enumerate}

\begin{enumerate}
\def\labelenumi{\arabic{enumi}.}
\setcounter{enumi}{29}
\tightlist
\item
  设 M 是 A-模，N 是 M 的一个子模，它有一个既约准素分解 \(N = \bigcap_{i=1}^{n} Q_{i}\) ；设 \(p_{i}\) 是伴随于 \(Q_{i}\) 的素理想。
\end{enumerate}

\begin{enumerate}
\def\labelenumi{(\alph{enumi})}
\item
  证明：若 \(\mathfrak{b}\) 是 \(\mathbf{A}\) 的理想且不含于任何 \(\mathfrak{p}_i\) 之中，则 \(\mathbf{N}:\mathfrak{b} = \mathbf{N}\) （参见习题 17）。
\item
  若诸 \(Q_{i}\) 中每一个在 M 中都是强准素的，反之证明：若 N: b = N，则 b 不含于任何 \(p_{i}\) 之中。（注意，若 b ⊂ \(p_{i}\) 且 P = \(\bigcap_{j \neq i} Q_{j}\) ，则对某个适当的整数 r 有 \(b^rP \subset N\) ，并由此推出，若 N: b = N，则 P = N。）由此得出，这时存在 \(\beta \in b\) 使 N: ( \(\beta\) ) = N。
\end{enumerate}

¶ 31. (a) 设 A 是 Noether 环，m 是 A 的极大理想，p 是含于 m 的素理想；证明：若 A 对于 m-进拓扑的 Hausdorff 完备化 \(\hat{A}\) 是整环，则 p 的符号幂 \(p^{(n)}\) 组成的滤子基对于 A 上的 m-进拓扑趋于 0。（化归到 \(\mathbf{A}\) 是局部环的情形；设 \(q\) 是 \(\hat{\mathbf{A}}\) 的素理想，它在 \(\mathbf{A}\) 上的迹为 \(p\) ，且对所有 \(n > 0\) ，设 \(c_{n} \notin p\) 使得 \(c_{n} p^{(n)} \subset p^{n}\) ；证明 \(c_{n} p^{(n)} \hat{\mathbf{A}} \subset q^{(n)}\) 。用习题 29 (e) 及第 III 章，§ 2，第 7 节，命题 8。）

\begin{enumerate}
\def\labelenumi{(\alph{enumi})}
\setcounter{enumi}{1}
\item
  设 A 是 Noether 环，n 是 A 的理想，p 是 \(\operatorname{Ass}(A/n)\) 中的极小素理想，\(m_{0}\) 与 m 是 A 的两个含 p 的极大理想；设 \(A(n)\) 、\(A(m_{0})\) 、\(A(m)\) 分别表示 A 对于 n-进、\(m_{0}\) -进、m-进拓扑的 Hausdorff 完备化。设 z 是 \(A(n)\) 的元素，它在 \(A(m_{0})\) 中的典范像为零。证明若 \(A(m)\) 是整环，则 z 在 \(A(m)\) 中的典范像也为零。（取 z 为 A 的元素的序列 \((x_{n})\) 的极限，使得 \(x_{n}-x_{m}\in n^{m}\) 对 \(n\geqslant m\) 成立，且 \(x_{n}\in m_{0}^{n}\) ；推出 \(x_{m}\) 属于 \(n^{m}\) 对于 \(m_{0}\) -进拓扑的闭包；用习题 8 断定 \(x_{m}\in p^{(m)}\) 对所有 m 成立，并借助 (a) 完成。）
\item
  设 A 是 Zariski 环，r 是 A 的定义理想；设 \(\operatorname{Spec}(A)\) 是连通的，且对 A 的每个极大理想 m，A 对于 m-进拓扑的 Hausdorff 完备化 \(A(m)\) 是整环。证明 A 对于 r-进拓扑的完备化 \(\hat{A}\) 是整环。（证明对 A 的每个极大理想 m，典范同态 \(\hat{A} \to A(m)\) 是单射。为此，可设 r 是有限多个素理想的交，这些素理想在 \(\operatorname{Ass}(A/r)\) 中极小，设 \(r = p_{1} \cap p_{2} \cap \cdots \cap p_{s}\) ；证明关于 \(\operatorname{Spec}(A)\) 的假设蕴涵对每个 i 存在极大理想 \(m_{i}\) 含 \((p_{1} \cap \cdots \cap p_{i-1}) + p_{i}\) 。用 (b) 证明：若 \(z \in \hat{A}\) 在 \(A(m)\) 中的典范像为零，且例如 \(p_{1} \subset m\) ，则 z 在每个 \(A(m_{i})\) 中的典范像为零，再断定 z 在 \(A(m')\) 中的典范像为零对 A 的每个极大理想 \(m'\) 成立。）
\end{enumerate}

\begin{enumerate}
\def\labelenumi{\arabic{enumi}.}
\setcounter{enumi}{31}
\tightlist
\item
  若环 \(\mathbf{A}\) 是局部环且其极大理想 \(\mathfrak{m}\) 是幂零理想 (*)，则称它为准素环。\(\mathbf{A}\) 中含于 \(\mathfrak{m}\) 的每个理想这时都是 \(\mathfrak{m}\) -准素的于 \(\mathbf{A}\) 之中。
\end{enumerate}

\begin{enumerate}
\def\labelenumi{(\alph{enumi})}
\item
  设 A 是强 Lasker 准素环，于是特别地其极大理想 m 是幂零的。证明：若 a, b 是含于 m 的两个理想，则必有 a: b ≠ a（用习题 30 (b)）。
\item
  准素 Noether 环是 Artin 环。由此推出：对 Noether 环 A 上的每个有限生成模 M 及 M 的每个 p-准素子模 N，存在整数 m，使得 M/N 的子模的每个严格递减序列 \((\mathrm{M}_{i}/\mathrm{N})_{0\leqslant i\leqslant k}\) （其中诸 \(M_{i}\) 是 p-准素的）长度 \(k\leqslant m\) （考虑 \(A_{p}\) -模 \(M_{p}\) ，并注意 \(M_{p}/N_{p}\) 被 \(pA_{p}\) 的一个幂所零化）。
\item
  证明一个 Artin 环 A，若其中每个素理想 \(\neq\{0\}\) 都是极大的，则它是有限多个准素环的直合成。
\end{enumerate}

\begin{enumerate}
\def\labelenumi{\arabic{enumi}.}
\setcounter{enumi}{32}
\tightlist
\item
  设 A 是交换环。若在 A/a 中零因子的集是一个理想 p/a，则称 A 的理想 a 为素性的；这时理想 p 是素的。
\end{enumerate}

\begin{enumerate}
\def\labelenumi{(\alph{enumi})}
\item
  证明每个准素理想与每个不可约理想都是素性的。
\item
  设 \(\mathfrak{a}\) 是 \(\mathbf{A}\) 的理想，使得理想 \(\mathfrak{b}\) （满足 \(\mathfrak{a}: \mathfrak{b} \neq \mathfrak{a}\) 者）组成的集有最大元 \(\mathfrak{p}\) ；证明 \(\mathfrak{p}\) 是素的，且 \(\mathfrak{a}\) 是素性的。* 给出一个素性理想 \(\mathfrak{a}\) 的例子，对它上述条件不满足（参见 § 1，习题 2）。*
\item
  在 Noether 环 A 中刻画素性理想，并给出一个非准素的素性理想的例子（参见 § 1，习题 11 (b)）。
\end{enumerate}

\begin{enumerate}
\def\labelenumi{\arabic{enumi}.}
\setcounter{enumi}{33}
\tightlist
\item
  在交换环 A 中，若对 A 的每对理想 b, c，关系 \(b \cap c \subset a\) 蕴涵 \(b \subset a\) 或 \(c \subset a\) ，则称理想 a 为拟素的；每个素理想都是拟素的，每个拟素理想都是不可约的。若中国剩余定理在 A 中成立，证明每个不可约理想都是拟素的（《代数》，第 VI 章，§1，习题 25）。
\end{enumerate}

* 35. 设 K 是交换域，B 是以 R 为阶群的赋值环，由``形式幂级数'' \(\sum_{x} c_{x} T^{x}\) 组成，其中 \(c_{x} \in \mathbf{K}\) ，\(x \in \mathbf{R}_{+}\) ，且诸 \(x\) 中使 \(c_{x} \neq 0\) 者组成的集是良序的；§1 习题 2 中定义的环 A 是 B 的子环，且 B 是平坦 A-模。若 C 是 §1 习题 2 中定义的 A-模，证明 \(\operatorname{Ass}_{\mathbf{B}}(\mathbf{C} \otimes_{\mathbf{A}} \mathbf{B}) \neq \varnothing\) ，尽管 \(\operatorname{Ass}_{\mathbf{A}}(\mathbf{C}) = \varnothing\) （考虑 C \(\otimes_{\mathbf{A}} \mathbf{B} = \mathbf{B}/\mathfrak{a}\mathbf{B}\) 中由 B 的一个赋值为 \(\alpha\) 的元素的典范像所得到的元素）。*

